\section{Errors in Numerical Methods}

A numerical solution is generally an approximation to an exact mathematical
quantity. Consequently, the analysis of errors associated with numerical
approximation is a fundamental component of numerical analysis
\cite{burden2016numerical,chapra2021numerical}.

Let $x$ denote the exact value of a quantity and $\tilde{x}$ a numerical
approximation. The absolute error may be defined as

\begin{equation}
    E_a = |x-\tilde{x}|,
\end{equation}

while the relative error can be expressed as

\begin{equation}
    E_r =
    \frac{|x-\tilde{x}|}{|x|},
\end{equation}

provided that $x \neq 0$.

Relative error is particularly useful because it relates the magnitude of the
approximation error to the magnitude of the quantity being approximated,
providing a scale-independent measure of accuracy
\cite{chapra2021numerical}.

In many practical numerical problems, however, the exact value is not known.
In such cases, stopping criteria based on successive approximations may be
used.

A common criterion consists of terminating the iterative procedure when

\begin{equation}
    |x_{n+1}-x_n| < \varepsilon,
\end{equation}

where $\varepsilon$ represents a predetermined tolerance.

The selection of the tolerance directly affects the computational behavior of
the algorithm. A relatively large tolerance may result in premature
termination, while an excessively small tolerance may increase the number of
iterations without necessarily producing a proportional improvement in
accuracy.

This limitation becomes especially important in finite-precision computation,
where the accuracy obtainable from an algorithm is ultimately constrained by
the arithmetic used by the machine
\cite{higham2002accuracy}.